% !TeX root = ../../Book5.tex
% 纲要标题：### 6.6 钩长公式与分支规则的应用
% 纲要位置：工作记录/计划纲要.md，第 3942 行。
% * 钩长公式的应用
% * 限制、诱导与迭代分支
% * 标准表与具体维数计算
% * 应用第6.5节已经证明的钩长公式与分支规则，系统计算低阶、钩形和两行分拆的表示维数，并核验分支两侧的维数恒等式；正式证明特征零下共轭分拆对应符号扭曲，以删格路径解释多步限制的重数
% ---
\section{钩长公式与分支规则的应用}
\label{b5:sec:0606}

钩长乘积、标准表递推和表示的分支分解分别从图形、组合和群作用三方面描述同一维数；可移格对应的递推项将把这三种描述联系起来，并给出低维计算中的相互检验。

\subsection{钩长公式的应用}
\label{b5:subsec:0606:01}

考虑 \(\lambda=(3,2,1)\)。按各格原位置排列，钩长为
\[
	\begin{array}{ccc}
		5&3&1\\
		3&1&\\
		1&&
	\end{array}.
\]
因此钩长乘积为 \(45\)，在特征零下，对应的 \(S_6\)-不可约表示维数为 \(6!/45=16\)。钩长计算中，每个格子本身只数一次；不能把同行格数与同列格数直接相加而忘记减去重复的格子。

\ProseParagraphBreak

转置图形保持所有钩长的多重集合，所以 \(f^\lambda=f^{\lambda'}\)。例如 \((4,1,1)\) 与 \((3,1,1,1)\) 的维数都为 \(10\)。在特征零下，转置图形还对应一个更强的表示论关系：符号扭曲把行求和与列交错求和互换。

\begin{proposition}\label{b5:prop:specht-conjugate-sign}
	设 \(k\) 为特征零域，\(n\) 为非负整数，\(\lambda\vdash n\)，\(\lambda'\) 为其共轭分拆，\(\operatorname{sgn}\) 为 \(S_n\) 的一维符号表示。在张量积上取对角左 \(S_n\) 作用，则有左 \(k[S_n]\)-模同构
	\[
		S_k^{\lambda'}\cong S_k^\lambda\otimes_k\operatorname{sgn}.
	\]
\end{proposition}
\begin{proof}
	令 \(A=k[S_n]\)。定义 \(k\)-线性映射 \(\alpha:A\to A\)，使 \(\alpha(g)=\operatorname{sgn}(g)g\) 对每个 \(g\in S_n\) 成立。符号的乘法性给出 \(\alpha(gh)=\alpha(g)\alpha(h)\)，而 \(\alpha^2=1\)，故 \(\alpha\) 是代数自同构。

	取形状 \(\lambda\) 的表 \(t\)，将它转置为形状 \(\lambda'\) 的表 \(t'\)。行群与列群互换，即 \(R_{t'}=C_t\)、\(C_{t'}=R_t\)，所以
	\[
		\alpha(a_t)=b_{t'},\qquad \alpha(b_t)=a_{t'},\qquad \alpha(c_t)=a_{t'}b_{t'},
	\]
	其中仍取 \(c_t=b_ta_t\)。由\cref{b5:prop:young-ideal-model}，\(S_k^\lambda\cong Ac_t\)。选符号表示的基向量 \(\eta\)，则
	\[
		Ac_t\otimes_k\operatorname{sgn}\longrightarrow Aa_{t'}b_{t'}, \qquad u\otimes\eta\longmapsto\alpha(u)
	\]
	是线性同构，因为 \(\alpha(Ac_t)=A\alpha(c_t)\)。它还满足
	\[
		\alpha\bigl(\operatorname{sgn}(g)gu\bigr) =g\alpha(u),
	\]
	所以与左 \(S_n\) 作用相容。

	还需把 \(Aa_{t'}b_{t'}\) 与 Specht 模所用的相反乘法次序联系起来。记 \(a'=a_{t'}\)、\(b'=b_{t'}\)，右乘 \(a'\) 给出左 \(A\)-模同态
	\[
		Aa'b'\longrightarrow Ab'a',\qquad u\longmapsto ua',
	\]
	因为 \((za'b')a'=(za')b'a'\) 属于目标左理想。此映射非零：由\cref{b5:prop:young-quasi-idempotent}，若 \(q=n!/f^\lambda\)，则
	\[
		(a'b')^2=\alpha(c_t^2)=q\,a'b'.
	\]
	这里 \(q\ne0\)，而 \(a'b'=\alpha(c_t)\ne0\)，后者也可由 \(c_t\) 的单位元系数为 \(1\) 看出。若 \(a'b'a'=0\)，则 \((a'b')^2=a'b'a'b'=0\)，与上式矛盾。因此右乘 \(a'\) 在生成元 \(a'b'\) 上的值非零。

	由\cref{b5:prop:specht-char-zero-simple}，\(Ac_t\) 与 \(Ab'a'\cong S_k^{\lambda'}\) 都是简单模；对第一个模作符号扭曲只把每个群元素的作用乘以非零标量，不改变其不变子空间，所以 \(Aa'b'\) 也简单。两个简单模之间的非零同态，核为零、像为整个目标，故上述右乘映射是同构。合成这两个同构即得结论。
\end{proof}

例如取特征零域 \(k\)，用 \(v_i\) 表示形状 \((3,1)\) 中第二行只有 \(i\) 的行表类。若一个表的第一列上、下两格分别为 \(a,b\)，它的列交错元就是 \(v_b-v_a\)。这些坐标差生成四点置换模的增广核
\[
	W=\left\{\sum_{i=1}^4 a_i v_i:\sum_{i=1}^4a_i=0\right\},
\]
故 \(S_k^{(3,1)}\cong W\)。这正是\subsecref{b5:subsec:0204:02}中的三维标准表示；命题于是给出 \(S_k^{(2,1,1)}\cong W\otimes_k\operatorname{sgn}\)。两个模的维数同为 \(3\)，它们之间的关系由符号扭曲具体实现。

\subsection{限制、诱导与迭代分支}
\label{b5:subsec:0606:02}

\((3,2,1)\) 有三个可移格，故特征零下
\[
	\operatorname{Res}_{S_5}^{S_6}S^{(3,2,1)} \cong S^{(2,2,1)}\oplus S^{(3,1,1)}\oplus S^{(3,2)}.
\]
三项维数为 \(5,6,5\)，合计 \(16\)。反过来，从 \((3,2)\) 添格得到
\[
	\operatorname{Ind}_{S_5}^{S_6}S^{(3,2)} \cong S^{(4,2)}\oplus S^{(3,3)}\oplus S^{(3,2,1)},
\]
维数为 \(9+5+16=30=6\cdot5\)，与诱导模维数相符。

\begin{corollary}\label{b5:cor:iterated-specht-branching}
	设 \(k\) 为特征零域，\(m,n\) 为满足 \(0\leq m\leq n\) 的整数，\(\lambda\vdash n\)、\(\mu\vdash m\)，\(S_k^\lambda,S_k^\mu\) 为相应 Specht 模，\(S_m\le S_n\) 为固定 \(m+1,\ldots,n\) 的子群。\(S_k^\mu\) 在 \(\operatorname{Res}_{S_m}^{S_n}S_k^\lambda\) 中的重数等于从 \(\mu\) 出发每次添一格到达 \(\lambda\) 的图形链条数。若 \([\mu]\subseteq[\lambda]\)，这个数也等于在差图 \([\lambda]\setminus[\mu]\) 中填入 \(m+1,\ldots,n\)，使行列均递增的表的个数；若没有该包含关系，重数为零。
\end{corollary}
\begin{proof}
	连续应用\cref{b5:thm:specht-branching}，每次删格所给成分重数均为 \(1\)，所以最终重数恰为经过 \(\mu\) 的删格链条数。反向读取链条，在从大小 \(a-1\) 到大小 \(a\) 新添的格子写 \(a\)。Young 图只能向右或向下扩张，所以每个新数字都大于它左边、上边已经填入的数字，得到所述表。

	反之，给定这种填法，依次保留 \([\mu]\) 及数字不超过 \(a\) 的格子。行列递增保证所得格子集向左、向上封闭，所以它仍是 Young 图；这恢复了原来的添格链条。两种操作互逆。
\end{proof}

差图称为\term[xie-Young-tu]{斜 Young 图}[skew Young diagram]。标准填数给出迭代分支重数的组合计算。

\begin{example}\label{b5:exm:iterated-branch-321}
	设 \(k\) 为特征零域，把 \(S_k^{(3,2,1)}\) 从 \(S_6\) 限制到固定数字 \(5,6\) 的子群 \(S_4\)。第一步得到形状 \((2,2,1),(3,1,1),(3,2)\)。第二步中，\((2,2,1)\) 删格得到 \((2,1,1)\) 与 \((2,2)\)，\((3,1,1)\) 得到 \((2,1,1)\) 与 \((3,1)\)，\((3,2)\) 得到 \((2,2)\) 与 \((3,1)\)。因此
	\[
		\operatorname{Res}_{S_4}^{S_6}S_k^{(3,2,1)} \cong \bigl(S_k^{(3,1)}\bigr)^{\oplus2} \oplus\bigl(S_k^{(2,2)}\bigr)^{\oplus2} \oplus\bigl(S_k^{(2,1,1)}\bigr)^{\oplus2}.
	\]
	维数为 \(2\cdot3+2\cdot2+2\cdot3=16\)。

	以最终形状 \(\mu=(3,1)\) 为例，两条删格路径是
	\[
		\begin{aligned}
			(3,2,1)&\longrightarrow(3,2)\longrightarrow(3,1),\\
			(3,2,1)&\longrightarrow(3,1,1)\longrightarrow(3,1).
		\end{aligned}
	\]
	差图 \((3,2,1)/(3,1)\) 的两个格子位于 \((2,2)\) 与 \((3,1)\)，它们既不同行也不同列，故没有行列递增所要求的先后限制。第一条路径先删 \((3,1)\)、再删 \((2,2)\)，因而在反向添格时分别在 \((2,2)\) 填 \(5\)、\((3,1)\) 填 \(6\)；第二条路径则在 \((3,1)\) 填 \(5\)、\((2,2)\) 填 \(6\)。这就是该斜图的两种标准填法，恰好对应 \(S_k^{(3,1)}\) 的重数 \(2\)。一次分支的每个成分只出现一次，但不同的连续分支路径可以在最后到达同一形状，使其重数累加。
\end{example}

\subsection{标准表与具体维数计算}
\label{b5:subsec:0606:03}

标准基还能直接给出生成元的矩阵。对 \(S^{(2,1)}\)，取
\[
	u=e_{\begin{smallmatrix}1&2\\3&\end{smallmatrix}}=v_3-v_1,
	\qquad
	v=e_{\begin{smallmatrix}1&3\\2&\end{smallmatrix}}=v_2-v_1.
\]
置换 \(s_1=(12)\)、\(s_2=(23)\) 满足
\[
	s_1u=u-v,\quad s_1v=-v,\qquad s_2u=v,\quad s_2v=u.
\]
因此在有序基 \(u,v\) 中，
\[
	\rho(s_1)=
	\begin{pmatrix}1&0\\-1&-1\end{pmatrix},
	\qquad \rho(s_2)=
	\begin{pmatrix}0&1\\1&0\end{pmatrix}.
\]
这两个整矩阵直接满足 \(\rho(s_i)^2=I\) 与 \(\rho(s_1)\rho(s_2)\rho(s_1)= \rho(s_2)\rho(s_1)\rho(s_2)\)。其迹在单位元、换位、三轮换上分别为 \(2,0,-1\)，与第一编的 \(S_3\) 特征标表一致。即使换到特征 \(3\)，这些矩阵仍定义一个表示，只是它会出现\cref{b5:exm:specht-mod-three}中的不变直线。

\ProseParagraphBreak

再看特征零域 \(k\) 上的 \(S_k^{(2,2)}\)。它有两个标准表
\[
	t_1=\begin{array}{cc}1&2\\3&4\end{array},
	\qquad
	t_2=\begin{array}{cc}1&3\\2&4\end{array}.
\]
令 \(x,y,z\) 分别表示无序配对 \(12|34,13|24,14|23\)，得到三维置换模 \(P=kx\oplus ky\oplus kz\)。这里两对数字没有编号；但行表类中的第一行与第二行有编号，因而 \(\{12|34\}\ne\{34|12\}\)，且 \(M_k^{(2,2)}\) 有 \(\binom42=6\) 个基向量。忘掉行号的等变满射
\[
	\pi:M_k^{(2,2)}\longrightarrow P
\]
把两行互换的两个表类送到同一配对，例如
\[
	\pi(\{12|34\})=\pi(\{34|12\})=x.
\]
它的核恰由三个向量
\[
	\{12|34\}-\{34|12\},\quad \{13|24\}-\{24|13\},\quad \{14|23\}-\{23|14\}
\]
生成：在每对互换的表类中，属于核的系数之和必须为零，故其贡献恰是这两个表类之差的倍数。因此 \(P\cong M_k^{(2,2)}/\ker\pi\)，这是从六维模到三维模的进一步商映射。

\ProseParagraphBreak

展开 \(t_1\) 的列反对称和，得到
\[
	e_{t_1}=\{12|34\}-\{23|14\}-\{14|23\}+\{34|12\}.
\]
第一项与第四项都映到 \(x\)，中间两项都映到 \(z\)；对 \(t_2\) 作同样计算，便有
\[
	\pi(e_{t_1})=2(x-z),\qquad \pi(e_{t_2})=2(y-z).
\]
因 \(2\) 可逆，这两个像构成坐标和为零平面的基，因此上述映射限制为同构
\[
	S_k^{(2,2)}\cong W, \qquad W=\{ax+by+cz:a+b+c=0\}.
\]
在当前特征零假设下，\(3\) 也可逆，故三点置换模分解为
\[
	kx\oplus ky\oplus kz=k(x+y+z)\oplus W.
\]
其中 \(k(x+y+z)\) 是平凡子模，\(W\) 是它的不变补空间。这个模型还表明 Klein 四元子群在 \(W\) 上平凡作用。

\ProseParagraphBreak

同一个忘行号映射 \(\pi\) 可以在任意域上定义。若 \(\operatorname{char}k=2\)，上面的两个像都为零，因此 \(\pi|_{S_k^{(2,2)}}=0\)。失效的是这个特定映射对 Specht 模的识别；标准基定理仍给出 \(\dim_kS_k^{(2,2)}=2\)，而\cref{b5:prop:specht-22-char-two}已经证明它在特征二下仍绝对不可约。不能由 \(\pi\) 的限制变成零映射，推断原 Specht 模消失或变得可约。

\subsection{低阶分拆与分支维数恒等式}
\label{b5:subsec:0606:04}

特征零下，\(S_4\) 的五个不可约维数按分拆次序为
\[
	\begin{array}{c|ccccc}
		\lambda &(4)&(3,1)&(2,2)&(2,1,1)&(1^4)\\ \hline f^\lambda&1&3&2&3&1
	\end{array},
	\qquad 1^2+3^2+2^2+3^2+1^2=24.
\]
例如 \((2,2)\) 虽有两行，但只有第二行末格可移，故它限制到 \(S_3\) 后仍为一个 \(S^{(2,1)}\)；\((3,1)\) 有两个可移格，限制则为 \(S^{(2,1)}\oplus S^{(3)}\)。

\ProseParagraphBreak

对 \(S_5\)，按
\[
	(5),(4,1),(3,2),(3,1,1),(2,2,1),(2,1,1,1),(1^5)
\]
排列，维数为 \(1,4,5,6,5,4,1\)，平方和为 \(120\)。逐个分拆的钩长计算给出了各项维数，平方和则核对它们与正则表示的维数是否一致。
